\begin{technical}
{\Large\textbf{Micromanipulation and Mitochondrial Heteroplasmy}}\\[0.3em]

\techheader{Technical Challenges in Nuclear Transfer}

Maternal spindle transfer requires extracting the metaphase II spindle-chromosome complex without disrupting chromosome alignment. The spindle, a 15–20-μm birefringent structure visible under polarised light (Oosight™), must be aspirated with minimal surrounding cytoplasm. Tachibana et al. (2009) demonstrated this technique in rhesus macaques, achieving no detectable donor mtDNA carry-over and producing healthy, fertile offspring.

Critical technical parameters:
\begin{itemize}[leftmargin=*]
\item \textbf{Timing}: Within 2 hours of oocyte retrieval before spindle depolymerization.
\item \textbf{Pipette diameter}: 20 μm bevelled, approaching at 30° to minimise membrane deformation.
\item \textbf{Cytoplasmic volume}: <5 picolitres co-aspirated to keep mitochondrial carry-over below 2\%.
\item \textbf{Fusion method}: Electrofusion (1.0 kV/cm, 50 μs pulses) or HVJ-E (Sendai virus extract).
\end{itemize}

Pronuclear transfer exploits the 8–10-hour window post-fertilisation when pronuclei are visible but unfused. Both pronuclei must be extracted together within a karyoplast—a membrane-bound cytoplasmic package containing $\sim$5\% of oocyte volume. This preserves their relative positioning, which carries epigenetic information essential for proper development.

\techheader{Mitochondrial Carry-over and Detection}

Even stringent micromanipulation cannot eliminate all donor mitochondria. Reported carry-over levels in embryo reconstruction studies include:
\begin{itemize}[leftmargin=*]
\item PNT: typically <2\%, occasionally higher (Hyslop et al., 2016).
\item MST: approximately 1–2\% on average (Tachibana et al., 2013).
\item PB1T: often <0.5\% when optimised (reported in polar body transfer studies).
\end{itemize}

Deep sequencing can achieve $\sim$0.1\% sensitivity, with digital droplet PCR detecting heteroplasmy as low as $\sim$0.01\% in optimised assays.

\techheader{Heteroplasmy Dynamics in Development}

Low-level heteroplasmy behaves unpredictably during development due to:

\textbf{The Mitochondrial Bottleneck}: During oogenesis, mtDNA copy number drops from $\sim$100,000 in mature oocytes to $\sim$200 in primordial germ cells before clonal expansion. This bottleneck allows random genetic drift to dramatically shift heteroplasmy levels between generations.

\textbf{Tissue-Specific Segregation}: Post-mitotic tissues show divergent heteroplasmy patterns:
\begin{itemize}[leftmargin=*]
\item \textbf{Muscle}: Can amplify from 1\% to 50\% by adulthood.
\item \textbf{Blood}: Typically maintains stable levels.
\item \textbf{Brain}: Shows regional variation, with high-energy regions (substantia nigra) potentially enriching for wild-type mtDNA.
\end{itemize}


\techheader{Nuclear-Mitochondrial Compatibility}

The mitochondrial proteome comprises ${\sim}$1,500 proteins: 13 encoded by mtDNA, the remainder nuclear-encoded and imported. OXPHOS complexes require precise stoichiometry between nuclear and mitochondrial subunits.

Potential incompatibilities arise where mtDNA-encoded subunits must interact with nuclear-encoded partners (e.g., Complex I: 7 mitochondrial and 38 nuclear subunits). However, the Tachibana primate studies (2009) showed reassuring intraspecies compatibility.

\techref
{\footnotesize
Hyslop, L. A. \textit{et al.} (2016). Towards clinical application of pronuclear transfer. \textit{Nature} \textbf{534}, 383–386.\\
Tachibana, M. \textit{et al.} (2009). Mitochondrial gene replacement in primate offspring and embryonic stem cells. \textit{Nature} \textbf{461}, 367–372.\\
Tachibana, M. \textit{et al.} (2013). Towards germline gene therapy of inherited mitochondrial diseases. \textit{Nature} \textbf{493}, 627–631.
}
\end{technical}